\section*{Capítulo \ref{ch:b2:geological-hazards} --- Peligros geológicos y riesgo}

\begin{solution}{exo:b2:geological-hazards:1}
(a) $1 - e^{-30/100} = 0.259$. (b) $e^{-100/100} = e^{-1} = 0.368$: más de un siglo de cada tres
pasa sin ninguna. (c) $T = -50/\ln 0.95 = 975$ años.
\end{solution}

\begin{solution}{exo:b2:geological-hazards:2}
(a) Ciudad A: $(1/200) \times 5000 \times 0.02 = 0.5$ edificios por año; ciudad B: $(1/50) \times 800
\times 0.10 = 1.6$. (b) La ciudad B, con 6.25 veces menos edificios expuestos. (c) Su peligro es 4 veces
más frecuente y su vulnerabilidad 5 veces mayor: en conjunto un factor 20, que compensa con creces la
menor \omterm{def:b2:geological-hazards:risk}{exposición}.
\end{solution}

\begin{solution}{exo:b2:geological-hazards:3}
(a) $\sqrt{9.81 \times 4000} = \qty{198}{m.s^{-1}}$; $\sqrt{9.81 \times 200} = \qty{44.3}{m.s^{-1}}$;
$\sqrt{9.81 \times 10} = \qty{9.90}{m.s^{-1}}$. (b) $\num{1.5e6}/198.1 = \qty{7570}{s}$, \qty{2.10}{h}.
(c) $800/3.6 = \qty{222.2}{m.s^{-1}}$, y $h = c^2/g = 222.2^2/9.81 = \qty{5030}{m}$.
\end{solution}

\begin{solution}{exo:b2:geological-hazards:4}
(a) $\tan\ang{33}/\tan\ang{25} = 0.6494/0.4663 = 1.39$. (b) $F = 1$ cuando $\beta = \varphi' =
\ang{33}$. (c) Con $c' = 0$, $F = (\cos^2\beta - r_u)\tan\varphi'/(\sin\beta\cos\beta)$, y $\cos^2\beta
- r_u = \cos^2\beta\,(1 - \gamma_w/\gamma)$, así que $F = (1 - \gamma_w/\gamma)\tan\varphi'/\tan\beta$:
el agua reduce el rozamiento a la mitad. $(1 - 9.81/20) \times 1.393 = 0.510 \times 1.393 = 0.710$: el
talud rompe.
\end{solution}

\begin{solution}{exo:b2:geological-hazards:5}
(a) $696/50 = 13.9$ por año. (b) $e^{-30/26.2} = 0.318$, frente a 0.304 observado: cercano. El
valor observado es algo menor porque las réplicas y los dobletes agrupados añaden intervalos
cortos. (c) $1 - e^{-7/26.2} = 0.234$.
\end{solution}

\begin{solution}{exo:b2:geological-hazards:6}
(a) El término de magnitud es $-0.105 + 0.229 \times 0.5 = 0.0095$. Para $d = 0$, 10, 30 y
\qty{100}{km}, $r = 5.57$, 11.45, 30.51 y \qty{100.2}{km}, $\log_{10}\mathrm{PGA} = -0.571$,
$-0.814$, $-1.145$ y $-1.547$, y la PGA mediana es 0.269, 0.153, 0.0715 y \qty{0.0284}{g}.
(b) $0.153 \times 10^{0.251} = 0.153 \times 1.78 = \qty{0.272}{g}$. (c) Una desviación típica
a cada lado: de $0.153/1.70 = \qty{0.0900}{g}$ a $0.153 \times 1.70 = \qty{0.260}{g}$.
\end{solution}

\begin{solution}{exo:b2:geological-hazards:7}
(a) Ordenados: 420 ($T = 13/1 = 13$ años), 305 (6.5), 268 (4.33), 251 (3.25). (b) $305 + (420 -
305)\ln(10/6.5)/\ln(13/6.5) = 305 + 115 \times 0.621 = \qty{376}{m^3.s^{-1}}$. (c) El mayor máximo
se representa en 13 años; cualquier valor a 100 años extrapola muy lejos del registro, con una ley de
extremos cuya elección y cuyos parámetros apenas acotan los 12 puntos.
\end{solution}

\begin{solution}{exo:b2:geological-hazards:8}
(a) $0.6 \times 0.4 \times 0.5 \times 0.3 = 0.036$. (b) $P(\text{erupción}) = 0.6 \times 0.4 = 0.24$;
sin flujo en el valle, $0.24 - 0.036 = 0.204$. (c) $0.6 \times 0.7 \times 0.5 \times 0.3 =
0.063$.
\end{solution}

\begin{solution}{exo:b2:geological-hazards:9}
Aquí $\sin\beta\cos\beta = 0.4698$, $\cos^2\beta = 0.6710$ y $\tan\ang{30}/\tan\ang{35} =
0.8245$. (a) $F = c'/(\gamma z\sin\beta\cos\beta) + 0.8245 = 1$ da $z_c = 8/(19 \times 0.4698
\times 0.1755) = \qty{5.11}{m}$. (b) $8/(19 \times 2 \times 0.4698) + 0.8245 = 0.448 + 0.825 =
1.27$. (c) $r_u = 9.81 \times 0.6710/19 = 0.3465$; $F = [8 + 38 \times (0.6710 - 0.3465) \times
0.5774]/(38 \times 0.4698) = 15.12/17.85 = 0.847$: la capa saturada rompe.
\end{solution}

\begin{solution}{exo:b2:geological-hazards:10}
(a) $\Pr(\text{ninguno en } [0, s+t] \mid \text{ninguno en } [0, s]) = e^{-(s+t)/T}/e^{-s/T} = e^{-t/T}$.
(b) Los 400 años de calma no cambian nada: $1 - e^{-50/475} = 0.0999$, un \qty{10}{\percent}. (c)
Contradice la hipótesis de que la tasa no depende del tiempo transcurrido: tras el rebote, la
tensión debe acumularse de nuevo, así que la tasa es baja justo después de un terremoto y crece después. Una zona
sismogénica contiene muchas fallas en etapas distintas de sus ciclos; la suma de muchas secuencias independientes
de sucesos se comporta casi como un proceso de Poisson, y la tasa es todo lo que se sabe de cada una.
\end{solution}

\begin{solution}{exo:b2:geological-hazards:11}
(a) $A = 0.5 \times (4000/10)^{1/4} = 0.5 \times 4.47 = \qty{2.24}{m}$. (b) $\lambda = cT$:
$198.1 \times 1200 = \qty{238}{km}$ en el océano profundo, $9.90 \times 1200 = \qty{11.9}{km}$ en
\qty{10}{m} de agua. (c) La amplitud ya no es pequeña frente a la profundidad (\qty{2.2}{m} en
\qty{10}{m}), así que la onda se empina y puede romper; parte de la energía se refleja o se concentra por la
forma del fondo marino y de las bahías; en tierra, el remonte depende de la pendiente y de la rugosidad del
terreno.
\end{solution}

\begin{solution}{exo:b2:geological-hazards:12}
(a) $10/5.8 + 4 = 1.72 + 4 = \qty{5.72}{s}$. (b) La onda S llega a \qty{50}{km} al cabo de $50/3.2 =
\qty{15.6}{s}$: $15.6 - 5.72 = \qty{9.90}{s}$ de aviso; a \qty{100}{km}, $31.3 - 5.72 =
\qty{25.5}{s}$. (c) $3.2 \times 5.72 = \qty{18.3}{km}$.
\end{solution}

\begin{solution}{pb:b2:geological-hazards:1}
Los logaritmos se conservan con cuatro decimales en los pasos intermedios.
\textbf{1.} $10^{4-5} = 0.1$, $10^{4-6} = 0.01$ y $10^{4-7} = 0.001$ por año.
\textbf{2.} 10, 100 y 1000 años.
\textbf{3.} $10^{4-7.5} = \num{3.16e-4}$ por año; $0.001 - \num{3.16e-4} = \num{6.84e-4}$ por año.
\textbf{4.} $0.1 - 0.01 = 0.09$ y $0.01 - 0.001 = 0.009$ por año.
\textbf{5.} $1 - e^{-0.01 \times 50} = 0.393$.
\textbf{6.} $1 - e^{-0.001 \times 50} = 0.0488$.
\textbf{7.} $r = \sqrt{10^2 + 5.57^2} = \qty{11.45}{km}$; $\log_{10} r = 1.0588$ y $-0.778 \times
1.0588 = -0.8237$.
\textbf{8.} $\mu = -0.105 + 0.229(M - 6) - 0.8237$: $-1.0432$, $-0.8142$ y $-0.6424$; medianas
0.0905, 0.153 y \qty{0.228}{g}.
\textbf{9.} $10^{-0.8142 + 0.230} = \qty{0.260}{g}$.
\textbf{10.} $z = (\log_{10} 0.3 + 0.8142)/0.230 = (-0.5229 + 0.8142)/0.230 = 1.266$, y $1 -
\Phi(1.266) = 0.103$.
\textbf{11.} $-0.105 + 0.229 \times 1.5 - 0.8237 = -0.5852$: \qty{0.260}{g}.
\textbf{12.} Las probabilidades de superar \qty{0.2}{g} ($\log_{10} 0.2 = -0.6990$) son 0.0672
($z = 1.497$), 0.308 ($z = 0.501$) y 0.597 ($z = -0.246$): $\lambda = 0.09 \times 0.0672 + 0.009
\times 0.308 + \num{6.84e-4} \times 0.597 = \num{6.05e-3} + \num{2.77e-3} + \num{4.08e-4} =
\num{9.23e-3}$ por año.
\textbf{13.} Para \qty{0.3}{g}: $z = 2.262$, 1.266 y 0.520, probabilidades 0.0118, 0.103 y
0.302; $\lambda = \num{1.07e-3} + \num{9.24e-4} + \num{2.06e-4} = \num{2.20e-3}$ por año.
\textbf{14.} Para \qty{0.4}{g} ($\log_{10} 0.4 = -0.3979$): $z = 2.805$, 1.810 y 1.063,
probabilidades 0.00251, 0.0352 y 0.144; $\lambda = \num{2.26e-4} + \num{3.17e-4} + \num{9.84e-5} =
\num{6.41e-4}$ por año.
\textbf{15.} $-\ln 0.9/50 = \num{2.107e-3}$ por año, un \omterm{def:b2:geological-hazards:return}{periodo de retorno} de 475 años.
\textbf{16.} Entre \qty{0.3}{g} (\num{2.20e-3}) y \qty{0.4}{g} (\num{6.41e-4}).
\textbf{17.} La pendiente es $\ln(6.41/22.0)/\ln(4/3) = -1.233/0.2877 = -4.29$, y $a = 0.3 \times
(2.107/2.20)^{-1/4.29} = 0.3 \times 1.0101 = \qty{0.303}{g}$.
\textbf{18.} En \qty{0.3}{g}, la franja 5--6 aporta el \qty{48.5}{\percent} de la tasa; la 6--7, el
\qty{42.1}{\percent}; y la 7--7.5, solo el \qty{9.4}{\percent}. Incluso el mayor terremoto tiene una mediana de
\qty{0.260}{g}, por debajo de \qty{0.3}{g}: el nivel solo se alcanza en la cola superior de la dispersión, y
los frecuentes terremotos moderados aportan la mayor parte.
\textbf{19.} $4/(19 \times 3 \times 0.4330) + \tan\ang{34}/\tan\ang{30} = 0.162 + 1.168 = 1.33$.
\textbf{20.} $r_u = 9.81 \times 0.75/19 = 0.387$; $F = [4 + 57 \times (0.75 - 0.387) \times
0.6745]/(57 \times 0.4330) = 17.95/24.68 = 0.727$: el talud saturado rompe.
\textbf{21.} Con el nivel freático a una fracción $m$ del espesor, $r_u = 0.387\,m$. $F = 1$ cuando
$0.75 - 0.387\,m = (24.68 - 4)/(57 \times 0.6745) = 0.538$: $m = 0.548$, un nivel freático al
\qty{55}{\percent} de la capa.
\textbf{22.} Al dividir $F = 1$ por $\gamma z$: $0.0702 + (0.75 - 0.4330\,k) \times 0.6745 = 0.4330 +
0.75\,k$, así que $k_c = (0.0702 + 0.5059 - 0.4330)/(0.2921 + 0.75) = 0.1431/1.042 = 0.137$.
\textbf{23.} $\log_{10} 0.137 = -0.8633 = -0.105 + 0.229(M - 6) - 0.8237$ da $M = 6.29$. Esos
terremotos se producen con una tasa de $10^{4 - 6.29} - 10^{-3.5} = \num{4.81e-3}$ por año, una vez cada 208 años aproximadamente,
y la PGA de 475 años de la pregunta~17 es más del doble de los \qty{0.137}{g} críticos: según este
análisis, el talud seco rompe con una sacudida muy inferior al movimiento de 475 años.
\textbf{24.} La \omterm{def:b2:geological-hazards:psha}{aceleración máxima del suelo} con una probabilidad del \qty{10}{\percent} de superarse en 50
años en la ciudad es de \textbf{\qty[text-series-to-math=true]{0.303}{g}}.
\end{solution}
